\documentclass[11pt]{article}
\usepackage[margin=25mm]{geometry}
\usepackage{fontspec}
\setmainfont{Latin Modern Roman}
\usepackage{amsmath,amssymb,mathtools}
\usepackage{unicode-math}
\setmathfont{Latin Modern Math}
\usepackage{xurl,hyperref}
\hypersetup{colorlinks=true,urlcolor=blue}
\setcounter{secnumdepth}{0}
\setlength{\parindent}{0pt}
\setlength{\parskip}{5pt}
\setlength{\emergencystretch}{3em}
\newcommand{\sref}[2]{\hyperlink{source:#1:#2}{[S#2]}}
\newcommand{\cataloguescope}[1]{\paragraph{Recorded target.}#1}
\begin{document}
\section{The Happy Ending Problem}\label{statement}

\subsection{Definitions and mathematical statement}

\paragraph{Shared notation.}$\mathbb N=\{1,2,\ldots\}$, and $\mathbb Z,\mathbb Q,\mathbb R,\mathbb C$ denote the integers, rationals, reals and complex numbers. Any other conventions are specified locally.


A finite set $P\subset\mathbb R^2$ is in general position if no three of its points are collinear. A subset $Q\subset P$ is in convex position if every point of $Q$ is a vertex of its convex hull. For $n\geq3$, let $g(n)$ be the least positive integer $N$ such that every general-position set $P\subset\mathbb R^2$ with $|P|=N$ contains a convex-position subset $Q\subset P$ with $|Q|=n$. The Erd\H{o}s--Szekeres theorem ensures that this minimum is finite.
\paragraph{Statement recorded in the supplied source.}
Determine $g(n)$ for every integer $n\geq3$. In particular, is
\[
g(n)=2^{n-2}+1\qquad\text{for every }n\geq3?
\] \sref{statement}{2}

\subsection{Short English statement}
How many points in general position guarantee a convex polygon with a prescribed number of vertices? Is the minimum always one more than the standard power-of-two lower bound?
\subsection{Sources}
\begin{description}
\item[\hypertarget{source:statement:1}{S1}] ULAM.AI and the UnsolvedMath Contributors, UnsolvedMath: A Curated Collection of Open Mathematics Problems, catalogue entry. CC BY 4.0. \url{https://huggingface.co/datasets/ulamai/UnsolvedMath}.
\item[\hypertarget{source:statement:2}{S2}] Andrew Suk, \emph{On the Erd\H{o}s--Szekeres convex polygon problem}, arXiv:1604.08657v2 (2016), Section 1 and footnote 1. \url{https://arxiv.org/pdf/1604.08657}.
\end{description}
\label{end:statement}
\end{document}
